令 $f_t[i]$ 为将前 $i$ 项恰好分成 $t$ 个非空连续段的最小费用。每段 $[l,r]$ 的费用为 $w(l,r)$，因此
\[
f_t[i]=\min_{t-1\le k<i}\{f_{t-1}[k]+w(k+1,i)\}.
\]
初始 $f_0[0]=0$，其余状态为 $\mathrm{inf}$。

令 $p_t[i]$ 为最优分割点，费用相同时取最大的一个。固定组数 $t$ 时，要求 $p_t[i]\le p_t[i+1]$。先计算中点 $i$ 的最优分割点 $k$，再递归计算左右两半：左边只枚举不超过 $k$ 的分割点，右边只枚举不小于 $k$ 的分割点。上一层必须已经全部算完，不能直接用这份实现计算依赖本层较早状态的一维 DP。

费用满足四边形不等式是常用的充分条件，即对 $a\le b\le c\le d$，
\[
w(a,c)+w(b,d)\le w(a,d)+w(b,c).
\]
加上上一层的费用后，不等式仍成立。若较长前缀的最优分割点更小，交换两个状态的分割点后，总费用不会增加。两个状态原本都已最优，所以交换后仍然最优，这与费用相同时取最大分割点的约定矛盾。这里不要求费用随区间扩大而不减。

代码用 $g$ 保存上一层，$f$ 保存本层。需要恢复分组时，另存每一层的 $p_t[i]$，从 $(m,n)$ 开始，每次由 $(t,i)$ 回到 $(t-1,p_t[i])$。

计算一层 DP 时，同一递归深度的枚举范围只在端点重叠，总枚举量为 $O(n)$。递归深度为 $O(\log n)$。若 $w(l,r)$ 可在 $O(1)$ 内计算，划分 $m$ 组的总时间复杂度为 $O(mn\log n)$。
